\section{Comparison with Alvarado, Ramon, and Wang}\label{app:arw}

Alvarado, Ramon, and Wang (ARW)~\cite{AlvaradoRamonWang2026} develop a
differential geometric approach to entropy optimization under general
multi-relational equality constraints. Our results treat joint upper
tails and entropy optimization with a fixed edge density and one fixed
subgraph density in the open oversaturated range. We compare ARW's
assumptions with those of \Cref{thm:joint-upper-tail,thm:fixed-densities}.
Then we present a construction suggesting that additional assumptions may be needed
to justify the finite-step conclusion for all global minimizers
in~\cite[Theorem~5.25]{AlvaradoRamonWang2026}.

\subsection{Comparison of assumptions}\label{app:arw-setting}

We state ARW's conditions for a single constraint $C(W)=u$, where
\[
  C:L^\infty_{\mathrm{sym}}([0,1]^2)\longrightarrow\RR,
  \qquad C(W):=\sum_{i=1}^q c_i t(F_i,W).
\]
Here $q\ge1$, $F_1,\dots,F_q$ are finite simple graphs, $c_i\in\RR$, and
$u\in\RR$ is the prescribed constraint value.
For an integer $k\ge1$, a \emph{$k$-step representation} consists
of a symmetric real matrix $A=(a_{ij})_{i,j=1}^k$ and a mass vector
$\pi=(\pi_1,\dots,\pi_k)$ with $\pi_i\ge0$ and $\sum_{i=1}^k\pi_i=1$.
Partitioning $[0,1]$ into consecutive intervals $S_i$ of length $\pi_i$ gives
\[
  W_{A,\pi}:=\sum_{i,j=1}^k a_{ij}\1_{S_i\times S_j}.
\]
The map $(A,\pi)\mapsto C(W_{A,\pi})$ is polynomial, and a representation
is \emph{feasible} when $C(W_{A,\pi})=u$.

With $\pi$ fixed, use the $k(k+1)/2$ independent entries $a_{ij}$, $i\le j$,
as coordinates on the space of symmetric matrices. The Jacobian is the row vector
\[
  D_A\bigl[C(W_{A,\pi})\bigr]:=
  \left(\frac{\partial C(W_{A,\pi})}{\partial a_{ij}}\right)_{1\le i\le j\le k}
  \in\RR^{1\times k(k+1)/2}.
\]
The value $u$ is \emph{regular}~\cite[Definition~5.5]{AlvaradoRamonWang2026}
if this row vector is nonzero at every feasible real-valued representation,
for every $k$.

For non-extremality, let $T_k$ be the set of constraint values attained
by $k$-step graphons:
\[
  T_k:=\left\{C(W_{A,\pi}):
    A=A^\top,\ 0\le a_{ij}\le1,\ \pi_i\ge0,\
    \sum_{i=1}^k\pi_i=1\right\}.
\]
A regular value $u$ is \emph{non-extremal} if
$u\in T_k\setminus\bigcup_{\ell\ge1}\partial T_\ell$ for some $k$.
This excludes boundaries of the attainable density sets at these block
sizes.

For a real-valued objective $I$, ARW's \emph{independence}
hypothesis~\cite[Definition~5.4 and Theorem~5.25]{AlvaradoRamonWang2026}
requires that for every $k\ge2$ there is a representation with $\pi_i>0$
for all $i$ at which both parameterized functions are analytic and
\[
  \operatorname{rank}
  \begin{pmatrix}
    D_A\bigl[C(W_{A,\pi})\bigr]\\
    D_A\bigl[I(W_{A,\pi})\bigr]
  \end{pmatrix}=2.
\]

Finally, a global minimizer $W$ of $I$ under $C(W)=u$ is \emph{isolated}
in cut distance if there is $\rho>0$ such that every feasible graphon $V$ satisfies
\[
  0<\delta_\cut(W,V)<\rho\quad\Longrightarrow\quad I(V)>I(W).
\]
Equivalently, no other global minimizing class lies within cut distance
$\rho$ of $W$.
Under the assumptions above, including isolation of every global minimizing
class, the result in~\cite[Theorem~5.25]{AlvaradoRamonWang2026} asserts that
all global entropy minimizers admit finite-step representations.

For the two constraint problems studied in this paper,
\Cref{thm:joint-upper-tail,thm:fixed-densities} do not assume ARW's analytic
independence, non-extremal regularity, or isolation.
For fixed edge and subgraph densities, we use local independence of their
derivatives at an optimizer $W$, expressed by the constraint
qualification in \Cref{def:constraint-qualification}:
\[
  \det\Gamma_W(U_1,U_2)\ne0
\]
for bounded symmetric $U_1,U_2$ supported in $\Omega_\gamma(W)$ for some
$\gamma>0$. We prove this for every optimizer in the open oversaturated
range (\Cref{prop:optimizer-nondegeneracy}), without asserting regularity
at all feasible representations. For joint upper tails,
\Cref{lem:upper-tail-multiplier} obtains multipliers by separation of convex
sets. We do not assume isolation of optimizing graphons;
our results do not exclude a sequence of distinct optimizing classes
converging to $W$ in cut distance.

\subsection{A construction testing the finite-step conclusion}\label{app:arw-entropy-counterexample}

Let
$W_{\mathrm h}(x,y):=\1_{\{x+y\le1\}}$ be the \emph{half-graphon}, and set
\[
  W_*:=\frac14+\frac12W_{\mathrm h}.
\]
Lov\'asz and Szegedy~\cite{LovaszSzegedy2011} show that the half-graphon is
finitely forcible among bounded real-valued symmetric kernels and that
affine transformations preserve this property. Thus $W_*$ is
\emph{finitely forcible} in this class:
there are finite simple graphs $F_1,\dots,F_\ell$ such that, for every
bounded symmetric kernel $U$,
\[
  t(F_i,U)=t(F_i,W_*)\quad(1\le i\le\ell)
  \quad\Longleftrightarrow\quad U\sim W_*.
\]
Here \emph{weak isomorphism}, denoted by $U\sim W$, means equality of
$t(F,U)$ and $t(F,W)$ for every finite simple graph $F$.
For graphons it is equivalent to having zero cut distance; write $[W]$ for
the equivalence class.

Define
\[
  q(W):=\sum_{i=1}^\ell\bigl(t(F_i,W)-t(F_i,W_*)\bigr)^2,
  \qquad
  g(W):=\bigl(e(W)-\tfrac14\bigr)q(W).
\]
Since
\[
  t(F\sqcup F',W)=t(F,W)t(F',W),\qquad t(K_1,W)=1,
\]
where $\sqcup$ denotes disjoint union and $K_1$ is the one-vertex graph,
$g$ is a finite real linear combination of graph densities.
We take $C=g$ and $u=0$, so the constraint is $g(W)=0$.
Up to a positive constant factor, ARW's entropy objective is
\[
  I(W):=\int_{[0,1]^2}I_0(W(x,y))\dd x\dd y,
  \qquad I_0(z):=z\log z+(1-z)\log(1-z),
\]
where $0\le z\le1$ and $0\log0:=0$.

\begin{prop}\label{prop:arw-entropy-counterexample}
The constraint $g(W)=0$ satisfies the independence and non-extremal
regularity hypotheses of~\cite[Theorem~5.25]{AlvaradoRamonWang2026},
including independence with the entropy objective.  Its global
entropy minimizers are exactly the two isolated classes $[1/4]$ and
$[W_*]$, and $W_*$ is not multipodal.
\end{prop}

\begin{proof}
The feasible graphons $W$ satisfy either $e(W)=1/4$ or $W\sim W_*$ by construction.
If $e(W)=1/4$, Jensen's inequality and the strict convexity of $I_0$ give
\[
  I(W)\ge I_0\bigl(e(W)\bigr)=I_0(1/4),
\]
with equality if and only if $W=1/4$ almost everywhere.
Weakly isomorphic graphons $U,W$ admit measure-preserving maps $\phi,\psi$
such that $U^\phi=W^\psi$ almost everywhere, so $I(U)=I(W)$.
Consequently, if $W\sim W_*$, then
\[
  I(W)=I(W_*)=\tfrac12 I_0(1/4)+\tfrac12 I_0(3/4)=I_0(1/4).
\]
Thus there are exactly two minimizing classes.
Since $W_*-\tfrac14=\tfrac12W_{\mathrm h}\ge0$, its cut norm equals its
integral. Relabeling preserves this integral and leaves the constant
graphon unchanged, so
\[
  \delta_\cut\bigl(\tfrac14,W_*\bigr)
  =\bigl\|W_*-\tfrac14\bigr\|_\cut
  =\frac12\int_{[0,1]^2}W_{\mathrm h}(x,y)\dd x\dd y=\frac14.
\]
Both minimizing classes are therefore isolated.
The degree distribution, the law of $d_W(X)$ for uniform $X\in[0,1]$,
is invariant under weak isomorphism because its moments are star densities.
For $W_*$,
\[
  d_{W_*}(x)=\frac34-\frac{x}{2},
\]
so this distribution is uniform on $[1/4,3/4]$.
A finite-step kernel has only finitely many degree values and
therefore cannot be weakly isomorphic to $W_*$.

Every feasible real-valued finite-step kernel $U$ consequently has
$q(U)>0$ and $e(U)=1/4$.  Writing
$Dg(U)[V]:=\left.\frac{\dd}{\dd t}g(U+tV)\right|_{t=0}$ for bounded
symmetric $V$, we obtain
\[
  Dg(U)[V]=q(U)\int_{[0,1]^2}V,
  \qquad Dg(U)[1]=q(U)>0.
\]
The constant direction is available at every step size, so $0$ is
regular.  Along constant graphons $c$ near $1/4$,
\[
  g(1/4)=0,\qquad
  \left.\frac{\dd}{\dd c}g(c)\right|_{c=1/4}=q(1/4)>0.
\]
Hence $0\in\operatorname{int}T_k$ for every $k$, with entries
restricted to $[0,1]$, and $0$ is non-extremal.

It remains to verify independence of $g$ and $I$: for every $k\ge2$,
we need a $k$-step representation at which the Jacobian of $(g,I)$ in
the entries of $A$ has rank $2$.  For $k=2$, take $\pi=(1/2,1/2)$ and
\[
  A_\delta:=\begin{pmatrix}1/4+\delta&1/4\\1/4&1/4-\delta\end{pmatrix},
  \qquad U_\delta:=W_{A_\delta,\pi},\qquad 0<\delta<1/4.
\]
With $\pi$ fixed, the Jacobian has two rows, for $g$ and $I$, and three
columns, for $a_{11},a_{12},a_{22}$.  We choose the columns
for $a_{11}$ and $a_{22}$.
Since $e(U_\delta)=1/4$, these columns form the matrix
\[
  \frac{\partial(g,I)}{\partial(a_{11},a_{22})}(A_\delta,\pi)
  =\frac14
  \begin{pmatrix}
    q(U_\delta)&q(U_\delta)\\
    I_0'(1/4+\delta)&I_0'(1/4-\delta)
  \end{pmatrix}.
\]
Its determinant is
\[
  \det\frac{\partial(g,I)}{\partial(a_{11},a_{22})}(A_\delta,\pi)
  =\frac{q(U_\delta)}{16}
     \bigl(I_0'(1/4-\delta)-I_0'(1/4+\delta)\bigr)\ne0,
\]
since $q(U_\delta)>0$ and $I_0'$ is strictly increasing.
Thus the full Jacobian has rank $2$.
Both functions are analytic at these interior entries.

For $k>2$, refine the two blocks into $k$ positive-mass subblocks.
Consider the two parameter directions that vary all refined entries
within the first or the second original diagonal block together.
The derivative of $(g,I)$ in these directions is the same nonsingular
$2\times2$ matrix as above, so the full refined Jacobian also has rank $2$.
All entries remain in $(0,1)$, preserving analyticity and proving
independence for every $k\ge2$.
\end{proof}

Thus the assumptions stated in~\cite[Theorem~5.25]{AlvaradoRamonWang2026}
do not exclude an isolated global minimizer with no finite-step representation.

This example suggests extending the regularity requirement to all global
optimizers, including those without finite-step representations.
Each density difference defining $q$ vanishes at $W_*$, so for every
bounded symmetric $V$,
\[
  q(W_*+tV)=O(t^2),\qquad g(W_*+tV)=O(t^2),\qquad Dg(W_*)[V]=0.
\]
Requiring a nonzero constraint derivative at every global optimizer would
therefore exclude this example; for several constraints, the corresponding
condition is surjectivity of their joint derivative.
A separate argument would be needed to establish that this stronger
condition suffices for the general finite-step conclusion.
